\documentclass[a4paper,11pt]{article}
\usepackage[utf8]{inputenc}
\usepackage[T1]{fontenc}
\usepackage[english]{babel}
\usepackage{amsmath,amssymb,amsfonts}
\usepackage{braket}
\usepackage{geometry}
\geometry{margin=2cm}
\usepackage{tikz}
\usetikzlibrary{arrows.meta,positioning,quotes}

\title{Geometric Representation of Amplitude Amplification\\in Grover's Algorithm}
\author{}
\date{}

\begin{document}

\maketitle

\section*{The Geometric Picture}

Grover's algorithm operates in a two-dimensional subspace of the full $2^n$-dimensional Hilbert space. The key insight is that the algorithm's state vector always remains in the plane spanned by two orthogonal basis vectors:

\begin{itemize}
  \item $\ket{s}$: the uniform superposition
    \[
    \ket{s} = \frac{1}{\sqrt{N}}\sum_{x=0}^{N-1}\ket{x},\qquad N=2^n.
    \]
  \item $\ket{w}$: the target (marked) state.
\end{itemize}

Define the orthogonal complement
\[
\ket{s'} = \frac{1}{\sqrt{N-1}}\sum_{x\neq w}\ket{x}.
\]
The full space is spanned by $\{\ket{w},\ket{s'}\}$, but we work in the $\{\ket{s},\ket{w}\}$ basis, which is \emph{not} orthogonal.

\section*{Initial State and Angle $\theta$}

The algorithm starts in $\ket{s}$. We can write
\[
\ket{s} = \sin\theta\ket{w} + \cos\theta\ket{s'},
\]
where
\[
\sin\theta = \frac{1}{\sqrt{N}},\qquad
\cos\theta = \sqrt{\frac{N-1}{N}},\qquad
\theta = \arcsin\left(\frac{1}{\sqrt{N}}\right).
\]

For large $N$, $\theta\approx\frac{1}{\sqrt{N}}$ is small.

\section*{Grover Iteration as Reflections}

A single Grover iteration $G = DO$ consists of two reflections:

\subsection*{1. Oracle Reflection $O$}

The oracle reflects about the hyperplane orthogonal to $\ket{w}$:
\[
O = I - 2\ket{w}\bra{w}.
\]
Geometrically, $O$ is a reflection across the line $\ket{s'}$ in our plane.

\subsection*{2. Diffusion Reflection $D$}

The diffusion operator reflects about $\ket{s}$:
\[
D = 2\ket{s}\bra{s} - I.
\]

\subsection*{Composition: Rotation by $2\theta$}

The product of two reflections is a rotation by twice the angle between their axes. Here, $DO$ rotates the state vector by angle $2\theta$ toward $\ket{w}$:
\[
\ket{s} \xrightarrow{G} \ket{\psi_1} \xrightarrow{G} \ket{\psi_2} \xrightarrow{G} \cdots \xrightarrow{G} \ket{\psi_R} \approx \ket{w}.
\]

After $R\approx\frac{\pi}{4}\sqrt{N}$ iterations, the state is nearly $\ket{w}$.

\section*{Geometric Diagram}

\begin{figure}[h]
\centering
\begin{tikzpicture}[
  scale=1.5,
  >=Stealth,
  every node/.style={font=\footnotesize},
  axis/.style={very thick, ->, black},
  vector/.style={thick, -Stealth, black},
  reflect/.style={dashed, gray},
  rotate/.style={blue, thick}
]

% Coordinate axes
\draw[axis] (0,0) -- (0,3.5) node[above] {$\ket{w}$};
\draw[axis] (0,0) -- (4.2,0) node[right] {$\ket{s'}$};

% Initial state |s>
\draw[vector, red] (0,0) -- (3.8,0.3) node[right] {$\ket{s}$};
\node[red,below right,font=\scriptsize] at (2.2,0.15) {$\theta$};

% Target state |w>
\draw[vector, green!70!black] (0,0) -- (0,3.2) node[above] {$\ket{w}$};

% |s> axis (diffusion reflection line)
\draw[dashed,blue] (0,0) -- (3.8,0.3) node[right] {$\ket{s}$};

% First oracle reflection (across |s'>)
\draw[dashed,gray] (0,0) -- (4.2,0);
\node[gray,above left,font=\tiny] at (2.1,0) {$O$};

% State after oracle
\draw[vector,orange] (3.8,0.3) -- (1.6,1.2) node[above right] {$\ket{\psi_1^-}$};

% Diffusion reflection line
\draw[reflect] (0,0) -- (3.8,0.3);
\node[blue,above left,font=\tiny] at (1.9,0.15) {$D$};

% State after diffusion (after 1 iteration)
\draw[vector,violet] (0,0) -- (2.8,1.1) node[above right] {$\ket{\psi_1}$};
\draw[rotate,double,blue] (3.8,0.3) to[bend left=20] (2.8,1.1)
  node[midway,right,font=\tiny] {$2\theta$};

% Second iteration (dashed for clarity)
\draw[vector,orange,dashed,opacity=0.7] (2.8,1.1) -- (1.0,1.9);
\draw[vector,violet,dashed,opacity=0.7] (0,0) -- (1.8,2.2);
\draw[rotate,double,blue,dashed,opacity=0.7] (2.8,1.1) to[bend left=20] (1.8,2.2);

% Optimal point
\draw[fill=red,opacity=0.3] (0.2,3.1) circle (0.08);
\node[red,above right,font=\scriptsize] at (0.3,3.1) {Optimal};

\end{tikzpicture}
\caption{Geometric picture of Grover's algorithm. Each iteration $G=DO$ rotates the state vector toward $\ket{w}$ by angle $2\theta$. The oracle $O$ reflects across $\ket{s'}$, while the diffusion operator $D$ reflects across $\ket{s}$. After $\sim\frac{\pi}{4}\sqrt{N}$ iterations, the state reaches the optimal point near $\ket{w}$.}
\label{fig:grover-geometry}
\end{figure}

\section*{Mathematical Derivation}

\subsection*{State After One Iteration}

Start with $\ket{s} = \sin\theta\ket{w} + \cos\theta\ket{s'}$.

\textbf{Oracle reflection:}
\[
O\ket{s}
= \sin\theta\,O\ket{w} + \cos\theta\,O\ket{s'}
= -\sin\theta\ket{w} + \cos\theta\ket{s'}
= \sin\theta\,(-\ket{w}) + \cos\theta\ket{s'},
\]
which makes angle $\pi-\theta$ with $\ket{w}$.

\textbf{Diffusion reflection:}
Since $D$ reflects about $\ket{s}$,
\[
D\ket{\psi^-} = 2(\bra{s}\psi^-)\ket{s} - \ket{\psi^-},
\]
where $\ket{\psi^-} = O\ket{s}$. Compute
\[
\bra{s}\psi^- = \sin\theta(-\sin\theta) + \cos\theta(\cos\theta)
             = \cos(2\theta).
\]
Thus
\[
D\ket{\psi^-}
= 2\cos(2\theta)\ket{s} - \ket{\psi^-}.
\]

In the $\{\ket{w},\ket{s'}\}$ basis, this becomes
\[
\ket{\psi_1} = \sin(3\theta)\ket{w} + \cos(3\theta)\ket{s'}.
\]
The angle with $\ket{s'}$ increased from $\theta$ to $3\theta$---a rotation by $2\theta$.

\subsection*{General Iteration}

After $k$ iterations,
\[
\ket{\psi_k} = \sin\bigl((2k+1)\theta\bigr)\ket{w}
             + \cos\bigl((2k+1)\theta\bigr)\ket{s'}.
\]
The probability of measuring $\ket{w}$ is
\[
P_k = \sin^2\bigl((2k+1)\theta\bigr).
\]
The optimal number of iterations is
\[
k_{\text{opt}} = \left\lfloor\frac{\pi}{4\sqrt{N}}\right\rfloor,
\]
which maximizes $P_k\approx 1$.

\section*{Amplitude Amplification Generalization}

The same geometric picture underlies \emph{amplitude amplification}, a generalization where the initial state $\ket{s}$ has arbitrary (small) overlap $\sin\theta$ with the ``good'' subspace $\mathcal{G}$ (not necessarily a single state). The same reflection operators amplify the projection onto $\mathcal{G}$ until it dominates.

This framework extends Grover's algorithm to quantum counting, collision finding, and many optimization problems.

\section*{Key Insights}

\begin{itemize}
  \item Grover's algorithm is a \emph{rotation} in a 2D subspace by fixed angle $2\theta$ per iteration.
  \item The oracle $O$ ``marks'' the target by phase flip (reflection).
  \item The diffusion $D$ amplifies via inversion about the mean (second reflection).
  \item Optimal iterations $\sim\frac{\pi}{4}\sqrt{N}$ yield $\sim 1-\mathcal{O}(1/N)$ success probability.
  \item Quadratic speedup $O(\sqrt{N})$ vs.\ classical $O(N)$.
\end{itemize}

\end{document}
